\centerline{\bf \Large  11 класс}

\problem{\textbf{1}} Пюрвя выписал на доску два числа $11^{\lg 121}$ и $10 \cdot 10^{\lg ^2 11}+11$. Определите, какое из чисел больше, и объясните почему?

\textbf{Решение:} Поскольку $11^{\lg 121}=\left(11^{\lg 11}\right)^2$, а $10 \cdot 10^{\lg ^2 11}+11=10 \cdot\left(10^{\lg 11}\right)^{\lg 11}+11=10 \cdot 11^{\lg 11}+11$, требуется сравнить числа $x^2$ и $10 x+11$, где $x=11^{\lg 11}$. Выражение $x^2-10 x-11=(x-11)(x+1)$ в точке $x=11^{\lg 11}>11$ принимает положительное значение, поэтому первое число больше.

\textbf{Критерии.} \\
\textit{7 баллов} $-$ Полностью верное решение\\
\textit{0-6 баллов} $-$ Все остальные продвижения\\
\textit{0 баллов} $-$ Любой ответ без объяснений\\

\problem{\textbf{2}} Решите задачу:

а) Арслан бежит из Элисты в Оргакин. Он тратит 35\% времени на путь в гору, 50\% — по равнине, а остальное время — с горы. Время его движения из Элисты в Оргакин и по той же дороге из Оргакина в Элисту одинаково, а его скорости в гору, с горы и по равнине постоянны, но различны. Во сколько раз быстрее Арслан бежит с горы, чем в гору?


б) Составьте и решите обратную задачу по схеме: ?; 50; $\dfrac{7}{3}$.

\textbf{Решение А:}

\textbf{ВАЖНО! Доказать равенство длины подъёма и длины спуска, если оно используется в решении.}

Пусть скорость Арслана с горы в $x \neq 1$ раз больше скорости в гору.

Подъём на первом участке (слева направо) занимал $0,35$ от времени пути, тогда спуск на первом участке (справа налево) занимает $0,35 / x$ - так как скорость во время спуска в $x$ раз больше, следовательно, время в $x$ раз меньше.

Равнина - второй участок.

Спуск на третьем участке (слева направо) занимал $0,15$ от времени пути, тогда подъём на третьем участке занимает $0,15 x$ - так как скорость во время подъёма в $x$ раз меньше, следовательно, время в $x$ раз больше.

$$
\begin{aligned}
0,35+0,5+0,15= & 0,35 / x+0,5+0,15 x \quad \\
& \Rightarrow \quad 0,15 x^2-(0,35+0,15) x+0,35=0 \\
& \Rightarrow(x-1)(x-7 / 3)=0 \quad \Rightarrow \quad x=7 / 3 .
\end{aligned}
$$

\textbf{Примерное условие Б:} Арслан бежит из Элисты в Оргакин. Он тратит некоторое количество времени на путь в гору, 50\% — по равнине, а остальное время — с горы. Время его движения из Элисты в Оргакин и по той же дороге из Оргакина в Элисту одинаково, а его скорости в гору, с горы и по равнине постоянны, но различны. Скорость в гору в $\dfrac{7}{3}$ раз больше скорости с горы. Какой процент от общего времени Арслан тратит на путь в гору из Элисты в Оргакин?

\textbf{Решение Б:} Пусть некоторый процент $100A$ он тратил на путь c горы, тогда если скорость Арслана с горы в $\dfrac{7}{3}$ раз больше скорости в гору, то из условия имеем:

$$
\begin{aligned}
0,5-A+A= & (0,5-A )/ \dfrac{7}{3}+A \dfrac{7}{3} \quad \Rightarrow \quad A=0,15 .
\end{aligned}
$$

Тогда на путь в гору тратит 50-15=35\%.

\textbf{Критерии.} \\
\textit{0-4 баллов} $-$ Решение пункта А\\
\textit{0-1 балла} $-$ Составление условия пункта Б \\
\textit{0-2 баллов} $-$ Решение пункта Б\\
\textit{0 баллов} $-$ Любой ответ без объяснений\\


\problem{\textbf{3}} Даны квадратные трёхчлены $f_1(x)=x^2+2 a_1 x+b_1, f_2(x)=x^2+2 a_2 x+b_2, f_3(x)=x^2+2 a_3 x+b_3, ...,f_{2025}(x)=x^2+2 a_{2025} x+b_{2025}$. Известно, что $a_1\cdot a_2\cdot a_3\cdot...\cdot a_{2025}=b_1\cdot b_2\cdot b_3\cdot ...b_{2025}>1$. Докажите, что хотя бы один из этих трёхчленов имеет два корня.

\textbf{Решение:} Предположим противное, тогда дискриминанты всех трёхчленов неположительны, то есть $a_k^2 \leqslant b_k(k=1,2,3...2025)$. Левые (а значит, и правые) части этих неравенств неотрицательны, поэтому их можно перемножить, получая $\left(a_1 a_2 a_3...a_{2025}\right)^2 \leqslant b_1 b_2 b_3...b_{2025}$, то есть $N^2 \leqslant N$, где $N=a_1 a_2 a_3...a_{2025}$. Но это противоречит неравенству $N>1$.

\textbf{Критерии.} \\
\textit{7 баллов} $-$ Полностью верное решение\\
\textit{0-6 баллов} $-$ Все остальные продвижения\\
\textit{0 баллов} $-$ Любой ответ без объяснений\\

\problem{\textbf{4}} 
Пусть $\omega_1$ и $\omega_2$ - две окружности, которые внутренне касаются друг друга в точке $A$. Пусть их общая касательная обозначена как $l$. Выберем произвольные точки $B$ и $C$ на $l$. Пусть две касательные из $B$ к $\omega_1$ и $\omega_2$ касаются их в точках $X, Y$, отличных от $A$, а касательные из $C$ к $\omega_1$ и $\omega_2$ касаются их в точках $W, Z$, также отличных от $A$. Докажите, что точки $X, Y, W, Z$ лежат на одной окружности.

\textbf{Решение:} Пусть $O_1, O_2$ - центры $\omega_1, \omega_2$. Пусть $P$ - произвольная точка на $l$. Пусть касательные из $P$ к $\omega_1, \omega_2$ касаются их в $Q, R$.

Я утверждаю, что $Q R$ проходит через фиксированную точку на $O_1 O_2$. Пусть $O_1 O_2 \cap Q R=S$. И пусть $P O_1 \cap A Q=D, P O_2 \cap A R=E, D E \cap O_1 O_2=T$. Так как $P A=P Q=P R$, $A O_1=Q O_1, A O_2=R O_2$, получаем $A D=D Q, A E=E R$. Следовательно, $A T=T S$. и $\angle A D P=\angle A E P=90^{\circ}$, значит, $A D E P$ - вписанный четырехугольник.

Также $P D \cdot P O_1=P A^2=P E \cdot P O_2$, так что $D E O_1 O_2$ - вписанный четырехугольник.
Следовательно, $T A^2=T D \cdot T E=T O_1 \cdot T O_2$ и и точка $T$ явно фиксирована. Таким образом, $S$ - фиксированная точка.

Возвращаясь к основной задаче, по утверждению $S X Y$ и $S Z W$ коллинеарны. и $S X \cdot S Y=$ $S A^2=S Z \cdot S W$, что означает, что $W X Y Z-$ вписанный четырехугольник. Готово.

\textbf{Критерии.} \\
\textit{7 баллов} $-$ Полностью верное решение\\
\textit{0-6 баллов} $-$ Все остальные продвижения\\
\textit{0 баллов} $-$ Любой ответ без объяснений\\

\problem{\textbf{5}} 
В течение нескольких дней учитель выбирал из класса нескольких учеников (хотя бы одного) на дежурство. В итоге оказалось, что в любые два дня среди дежуривших учеников есть ровно один человек, который дежурил в оба этих дня. Докажите, что они не могли дежурить больше дней, чем есть учеников в классе (то есть, если в классе 30 человек, то дежурство не могло длиться дольше 30 дней).

\textbf{Решение:} Контрпример --- в классе всего один ученик, и он может дежурить сколько угодно  дней.

\textbf{Критерии.} \\
\textit{7 баллов} $-$ Полностью верное решение\\
\textit{0-1 балл} $-$ Все остальные продвижения\\
\textit{0 баллов} $-$ Любой ответ без объяснений\\